% =====================================================================
% Section 0 — Abstract
%
% Written last, after Sections 4, 5 and 6 were complete.
% CONSTRAINT: every number here already appears in Section 4. No result is
% introduced that is not reported in Results.
%
%   92.356 +/- 1.022   Table 1 (tab:task_b)
%   99.561 +/- 0.354   Table 1 (tab:task_b)
%   98.371 +/- 0.319   Table 1 (tab:task_b)
%   1.19 points        Sec 4.6
%   4.22 M / 27.83 M   Table 4 (tab:efficiency)
%   61 vs 1            Sec 4.3
%   95.3 / 96.6 / 37.970  Sec 4.5
%   17 species, 532 test images, 3 seeds  Sec 4.1
% =====================================================================

\begin{abstract}
Macroscopic cross-section imaging is the most deployable route to conifer
wood identification, but it is unclear which modern architecture it calls
for, or how much of the reported accuracy reflects anatomical evidence. We
benchmark ResNet-50, EfficientNet-B3, MobileNetV3-Large, ConvNeXt-Tiny and
Swin-Tiny on 17 conifer species, using a partition grouped by origin key,
three seeds per configuration, and a single immutable evaluation registry.
Results are reported as mean and standard deviation across seeds; no
interval estimate is claimed.

ResNet-50 is not competitive, reaching $92.356 \pm 1.022$\,\% test accuracy
against $99.561 \pm 0.354$\,\% for ConvNeXt-Tiny, and producing 122 of the
204 errors made by all backbones and seeds combined. The remaining four
architectures are separated by margins comparable to their seed-to-seed
variation, so three seeds do not establish their ordering.
MobileNetV3-Large is the practically notable case, reaching
$98.371 \pm 0.319$\,\% with 4.22\,M parameters against 27.83\,M. Errors are
concentrated: in the strongest model every error is the same confusion
between \textit{Pinus krempfii} and \textit{Pinus dalatensis}, a pair
present in all five architectures and asymmetric at 61 occurrences against
one reversed. Two augmentation ablations were each trained with a single
seed and support no recommendation.

A control experiment qualifies these numbers. Models trained on
$16\times16$ or heavily blurred inputs retain 95.3\,\% and 96.6\,\% of
full-image accuracy, while a colour-histogram-only model reaches
37.970\,\%. Since anatomical structure does not survive those
transformations, the accuracies reported here describe performance on this
corpus rather than anatomical identification, and external validation
remains the necessary next step.
\end{abstract}

\begin{paperkeywords}
wood identification, conifer species, macroscopic cross-section, deep
learning benchmark, shortcut learning, dataset leakage
\end{paperkeywords}
